\documentclass[12pt, a4paper]{article}

\usepackage[utf8]{inputenc}
\usepackage[spanish]{babel}
\usepackage{geometry}
\geometry{
    a4paper,
    left=25mm,
    right=25mm,
    top=25mm,
    bottom=25mm
}
\usepackage{enumitem}
\usepackage{xcolor}
\usepackage{hyperref}
\usepackage{titlesec}
\usepackage{graphicx}
\usepackage{mdframed}
\usepackage{parskip} % Espaciado entre parrafos sin sangria

% Configuración de colores
\definecolor{primary}{RGB}{41, 128, 185}
\definecolor{darkgray}{RGB}{60, 60, 60}

% Configuración de enlaces
\hypersetup{
    colorlinks=true,
    linkcolor=primary,
    filecolor=magenta,      
    urlcolor=primary,
}

% Cajas para resaltados (Contexto, Pregunta, Resultado)
\newmdenv[
  leftline=true,
  rightline=false,
  topline=false,
  bottomline=false,
  linewidth=4pt,
  linecolor=primary,
  backgroundcolor=gray!10,
  innerleftmargin=10pt,
  innerrightmargin=10pt,
  innertopmargin=10pt,
  innerbottommargin=10pt
]{highlightbox}

\begin{document}

\begin{center}
    \vspace*{1cm}
    \Huge \textbf{Taller de Innovación Nestlé} \\
    \vspace{0.5cm}
    \LARGE \textbf{Propuesta de Valor AeroStock OS} \\
    \vspace{1cm}
    \rule{\linewidth}{0.5mm}
\end{center}

\vspace{0.5cm}
Este documento sintetiza la información operativa de Nestlé, las proyecciones tecnológicas y los entregables requeridos para la ronda final de pitch.

\section{Análisis del Problema y Expectativas}

\subsection*{El Problema Actual (Dolores / Pain Points)}
\begin{itemize}[label=\textcolor{red}{\textbullet}]
    \item \textbf{Inventario Manual e Inseguro:} El almacén cuenta con \textbf{24,000+ posiciones} a \textbf{12 metros de altura}. Al tener \textbf{$\sim$0\% de automatización}, la auditoría exige elevar operarios en canastas de montacargas, generando un alto riesgo de accidentes.
    \item \textbf{Trazabilidad Lenta y Desconectada:} El WMS se actualiza de forma manual post-ingreso. Esta falta de sincronización en tiempo real con SAP genera altas tasas de error y desajustes físicos.
    \item \textbf{Cuellos de Botella Constantes:} La dependencia intensiva de montacargas manuales crea bloqueos al final de la línea. Esto fuerza paros no planificados que reducen el \textbf{OEE (Overall Equipment Effectiveness)} de la planta 24/7.
\end{itemize}

\subsection*{El Desafío (Expectativas de la Solución)}
\begin{itemize}[label=\textcolor{green}{\textbullet}]
    \item \textbf{Capacidad:} Mover y auditar un flujo de \textbf{215 pallets por turno} en un tramo de \textbf{20 metros}.
    \item \textbf{Continuidad:} Operación \textbf{24/7} con \textbf{cero interrupciones} no planificadas.
    \item \textbf{Integración:} Sincronización instantánea de WMS y SAP.
    \item \textbf{KPIs Objetivo:} Cero incidentes de seguridad, minimización de tiempos de ciclo y maximización de la exactitud del inventario.
\end{itemize}

\section{Formato: Contexto + Pregunta + Resultado}

\begin{highlightbox}
\textbf{(Contexto con cuantificación):} \\
Actualmente en el almacén, auditar manualmente racks de hasta 12 metros de altura expone a los operarios a riesgos fatales y consume extensas jornadas. Esto perpetúa un nivel de automatización del $\sim$0\% y retrasa el flujo crítico de 215 pallets por turno (dependientes de montacargas manuales). Estos paros impactan gravemente el OEE y la capacidad productiva de la planta.
\end{highlightbox}

\begin{highlightbox}
\textbf{(Pregunta Reto):} \\
¿Cómo podríamos automatizar la búsqueda y validación de productos en posiciones de gran altura, de tal manera que reduzcamos drásticamente el tiempo de conteo, eliminemos al 100\% el riesgo humano y sincronicemos la trazabilidad en tiempo real con SAP y WMS?
\end{highlightbox}

\begin{highlightbox}
\textbf{(Resultado Esperado):} \\
Reducir los ciclos de inventario a rutinas diarias automatizadas que alcancen un 99.9\% de exactitud en los datos, erradicando por completo la necesidad de realizar trabajo humano en alturas peligrosas.
\end{highlightbox}

\section{Diseño de Propuesta de Valor}

\begin{itemize}
    \item \textbf{NUESTRO(S):} AeroStock OS (Sistema ASRS visual integrado con Drones).
    \item \textbf{AYUDA(N):} A los jefes de bodega y operarios logísticos de Nestlé.
    \item \textbf{QUE QUIEREN:} Mantener un flujo ininterrumpido de 215 pallets por turno (24/7) y tener control total sobre sus inventarios.
    \item \textbf{PARA (Analgésico):} Eliminar la accidentabilidad por caídas de altura (12m) y reducir los frustrantes cuellos de botella.
    \item \textbf{Y (Creador de Valor):} Aumentar la exactitud del inventario por encima del 99\% mediante auditorías visuales (Lectura LPN y detección de caducidades) sincronizadas instantáneamente.
    \item \textbf{A DIFERENCIA DE:} Los métodos tradicionales manuales, donde el personal debe trepar en canastas sin validación inmediata de datos.
\end{itemize}

\section{Cuantificación y Estimación de Mejora}

Proyecciones de impacto del prototipo AeroStock frente al proceso manual actual:

\begin{enumerate}
    \item \textbf{Eficiencia (Reducción de Tiempo):} Un operador tarda 2 a 3 minutos por cada pallet a gran altura. Un dron ASRS vuela y escanea en 2 a 4 segundos por pallet. Esto representa una \textbf{reducción del 85\% al 90\%} en tiempos de captura.
    \item \textbf{Seguridad (Cero Riesgo de Altura):} El 100\% de las posiciones superiores a 2 metros exigen maquinaria para elevar personas. El impacto es del \textbf{100\% de eliminación} de horas-hombre en trabajo riesgoso.
    \item \textbf{Calidad (Trazabilidad y Precisión):} La tasa de error humano ronda el 3\% y 5\%. La visión artificial garantiza una exactitud del \textbf{99.9\%} con alertas directas.
\end{enumerate}

\section{Estructura de Pitches (Estrategia: Demo en Vivo Integrada)}

\subsection*{Pitch Preliminar (1 minuto)}
\textit{``Hola, somos el equipo AeroStock. Actualmente en Nestlé operamos con más de 24,000 posiciones a 12 metros de altura y el control de inventario es 100\% manual. Esto obliga a personas a subir en canastas de montacargas, poniendo en riesgo sus vidas y causando cuellos de botella que frenan la planta 24/7.}

\textit{Para solucionarlo, creamos \textbf{AeroStock OS}, un sistema de visión artificial que usa drones autónomos para auditar la bodega. Escanea códigos LPN y detecta caducidades en tiempo real, validándolo con el WMS. Eliminamos el 100\% del riesgo en alturas y logramos un 99.9\% de exactitud. ¡Tenemos un prototipo físico funcional aquí mismo que les mostraremos en la ronda final! Gracias.''}

\subsection*{Pitch Final (3 minutos - Incluyendo Demo en Vivo)}

\textbf{Diapo 1: Contexto y Problema (30 seg)} \\
``¿Sabían que la falta de automatización logística afecta directamente el OEE de nuestra planta? Hoy controlamos 24,000 posiciones manualmente. Subir operarios a 12 metros de altura para auditar es un riesgo de vida latente y genera cuellos de botella que interrumpen nuestra producción 24/7. Nuestro registro es tardío y no habla con SAP en tiempo real.''

\textbf{Diapo 2: La Solución Teórica (30 seg)} \\
``La pregunta fue: ¿Cómo auditar a 12 metros en segundos sin arriesgar a nadie? Les presentamos \textbf{AeroStock OS}. Un cerebro de Inteligencia Artificial que recibe video de drones. Mientras el dron vuela orientándose con AprilTags, lee LPNs y clasifica el producto. Si hay un error, alerta al WMS. Cero intervención humana.''

\textbf{TRANSICIÓN A DEMO EN VIVO (Min. 1:00 a 2:15 - 75 seg)} \\
\textit{[Nota: Quitan las diapositivas y muestran la pantalla web de AeroStock (Flask). Alguien tiene el celular (Dron Principal) apuntando a las 3 cajas físicas en la mesa].} \\
``Pero no vinimos a contárselo, vinimos a demostrárselo en vivo. Tenemos aquí 3 cajas físicas simulando los Racks A1, A2 y A3. Esta cámara web de nuestro celular simula el dron en vuelo. Observen la pantalla: pasamos el `dron' por el primer Rack... [hacen el escaneo de la primera caja]. ¡Boom! Detecta el código LPN al instante. Ahora pasamos a este otro Rack... [escanean la caja con caducidad]. ¡Fíjense en la alerta roja de CADUCADO! El sistema bloquea el producto en tiempo real. Y si vemos el `Historial', todo queda registrado con ubicación y fecha exacta para auditoría.''

\textbf{Diapo 3: Propuesta de Valor y Resultados (30 seg)} \\
\textit{[Nota: Regresan rápido a las diapositivas]} \\
``Como acaban de ver, transformamos un proceso peligroso en una auditoría continua:
\begin{itemize}
    \item \textbf{Seguridad:} CERO riesgo en altura.
    \item \textbf{Eficiencia:} Reducimos el tiempo en más de un 85\%.
    \item \textbf{Trazabilidad:} Exactitud del 99.9\%.
\end{itemize}
Aseguramos que el flujo de 215 pallets no sufra bloqueos, protegiendo el OEE.''

\textbf{Diapo 4: Cierre (15 seg)} \\
``Con AeroStock OS, el almacén deja de ser un cuello de botella y se convierte en una ventaja competitiva ágil, rentable y segura. Muchas gracias.''

\end{document}
